\chapter{Compiling LLVM IR for Different CPU Architectures}

One of LLVM IR's key strengths is \textbf{write once, compile anywhere}. The same IR can target x86-64, ARM, RISC-V, or WebAssembly by changing the compilation flags.

\section{Target Triples}

A target triple specifies the architecture, vendor, and operating system:

\begin{lstlisting}[style=ts]
import { LLVMCodegen } from "@euriklis/llvm-ir";

const codegen = new LLVMCodegen("my_module");
codegen.targetTriple = "x86_64-unknown-linux-gnu";  // x86-64 Linux
codegen.dataLayout = "e-m:e-p270:32:32-i64:64-f80:128-n8:16:32:64-S128";
\end{lstlisting}

Common triples:
\begin{lstlisting}[style=ir]
x86_64-unknown-linux-gnu       # Linux x86-64
x86_64-apple-darwin            # macOS Intel
aarch64-unknown-linux-gnu      # Linux ARM64 (Raspberry Pi, AWS Graviton)
aarch64-apple-darwin           # macOS Apple Silicon (M1/M2/M3)
armv7-unknown-linux-gnueabihf  # ARM 32-bit (Raspberry Pi 3)
riscv64-unknown-linux-gnu      # RISC-V 64-bit
wasm32-unknown-unknown         # WebAssembly 32-bit
wasm64-unknown-unknown         # WebAssembly 64-bit
\end{lstlisting}

\section{Compiling with llc}

Once you have LLVM IR, compile it with \texttt{llc} (LLVM static compiler):

\begin{lstlisting}[style=ir]
# Generate x86-64 assembly
llc -march=x86-64 -filetype=asm -o output.s module.ll

# Generate ARM64 object file
llc -march=aarch64 -filetype=obj -o output.o module.ll

# Generate RISC-V binary
llc -march=riscv64 -filetype=obj -o output.o module.ll

# Generate WebAssembly
llc -march=wasm32 -filetype=obj -o output.wasm module.ll
\end{lstlisting}

Or use \texttt{clang} to link directly:
\begin{lstlisting}[style=ir]
clang -o program module.ll
./program
\end{lstlisting}

\section{Architecture-Specific Code with this Library}

The \texttt{@euriklis/llvm-ir} package provides architecture-specific code generators:

\textbf{x86-64 (default):}
\begin{lstlisting}[style=ts]
import { LLVMCodegen } from "@euriklis/llvm-ir";
const codegen = new LLVMCodegen("x86_module");
// Targets x86-64 by default
\end{lstlisting}

\textbf{ARM/AArch64:}
\begin{lstlisting}[style=ts]
import { ARMCodegen } from "@euriklis/llvm-ir";
const codegen = new ARMCodegen("arm_module", true);  // true = 64-bit
// Generates: target triple = "aarch64-unknown-linux-gnu"
\end{lstlisting}

\textbf{RISC-V:}
\begin{lstlisting}[style=ts]
import { RISCVCodegen } from "@euriklis/llvm-ir";
const codegen = new RISCVCodegen("riscv_module");
// Generates: target triple = "riscv64-unknown-elf"
\end{lstlisting}

\textbf{WebAssembly:}
\begin{lstlisting}[style=ts]
import { WASMCodegen } from "@euriklis/llvm-ir";
const codegen = new WASMCodegen("wasm_module");
// Generates: target triple = "wasm32-unknown-unknown"
\end{lstlisting}

All use the same API—only the code generator class changes.

\section{Data Layout Strings}

The data layout string specifies memory alignment, pointer sizes, and endianness:

\begin{lstlisting}[style=ir]
e-m:e-p270:32:32-i64:64-f80:128-n8:16:32:64-S128
\end{lstlisting}

Breakdown:
\begin{itemize}[noitemsep]
  \item \texttt{e}: Little-endian (Intel, ARM). \texttt{E} would be big-endian.
  \item \texttt{m:e}: ELF mangling (Linux).
  \item \texttt{p270:32:32}: Pointers in address space 270 are 32-bit with 32-bit alignment.
  \item \texttt{i64:64}: 64-bit integers have 64-bit alignment.
  \item \texttt{f80:128}: 80-bit floats (x87) have 128-bit alignment.
  \item \texttt{n8:16:32:64}: Native integer widths (CPU can operate on these efficiently).
  \item \texttt{S128}: Stack alignment is 128 bits (16 bytes).
\end{itemize}

Each architecture has its own data layout. The library sets these automatically when you use architecture-specific code generators.

\section{Cross-Compilation Example}

\textbf{Same TypeScript code, three architectures:}

\begin{lstlisting}[style=ts]
import {
  LLVMCodegen,
  ARMCodegen,
  WASMCodegen,
  LLVMFunction,
  BasicBlock,
  RetInst
} from "@euriklis/llvm-ir";

function createAddFunction(codegen) {
  const func = new LLVMFunction({
    name: "add",
    returnType: LLVMCodegen.types.i32
  });
  // ... (same function body for all architectures)
  codegen.getModule().addFunction(func);
}

// x86-64
const x86 = new LLVMCodegen("x86");
createAddFunction(x86);
console.log("=== x86-64 ===");
console.log(x86.toString());

// ARM64
const arm = new ARMCodegen("arm", true);
createAddFunction(arm);
console.log("=== ARM64 ===");
console.log(arm.toString());

// WebAssembly
const wasm = new WASMCodegen("wasm");
createAddFunction(wasm);
console.log("=== WASM ===");
console.log(wasm.toString());
\end{lstlisting}

Each will have a different target triple, but the function IR is identical. LLVM's backends handle architecture-specific optimizations (register allocation, instruction selection, calling conventions).

